\section{Trace identities and prime substitutions}
\label{sec:trace-calculus}

We now turn the least-depth coefficients into rules for replacing,
contracting, and deleting primes. The input is a coefficient system
from Section~\ref{cs:coefficient-systems}, with fixed support $S$,
characters $\eta,\psi$, expansion $g=\sum_{n\ge1}a(n)q^n$,
minimum depth $\delta$, and symbol $A(D)\in\F_2$, extended by zero
off its filter. The coefficient lower bounds and the rational leading
units proved there are essential hypotheses. A unit symbol has the
analytic interpretation of Proposition~\ref{prop:unit-symbol-detection}.

Put $F=\Q(E[2])$.  A prime outside $S$ is called split, simple, or a root
according as its action on $E[2]$ is the identity, a transposition, or an
element of order three; its weight is respectively $1$, $1/2$, or $0$.
The same terminology will be used for Galois elements, according to their
image in $\Gal(F/\Q)$.  A nonroot is an element of either of the first two
types.  A finite set of primes and its squarefree product will occasionally
be denoted by the same letter.  Thus $U/p$ means deletion of $p$ from $U$,
not deletion of a prime power from an arbitrary integer.

An oriented Frobenius label in an extension $L$ of $F$ means a
\emph{rational} arithmetic Frobenius lift in $\Gal(L/\Q)$, with
a chosen place and its residual image in $\Gal(F/\Q)$ retained.
At a simple prime this is the lift $g$ of the transposition, not
only the Frobenius $g^2$ of the residue-degree-two place of $F$.

For the odd source, fix the finite operator patterns, their weight bounds,
and all integer divisions required by the intended tests before choosing
an actual stage. The primes realizing those patterns may depend on the
stage. The lemma below distinguishes the sharp estimate on a prescribed
bounded range of squarefree weights from integrality at every Fourier
index. Both are needed to form the integral weight-two trace lattice. At a
chosen stage the detecting prime belongs to $S$, and the level,
coefficient field, and characteristic-zero lift are fixed. Recurrences
hold modulo the working precision before division.

\subsection{Normalized coefficient operators}

Let $U$ be a finite set of nonroot primes outside $S$.  Write
\[
 \mathsf U_p\!\left(\sum_n c(n)q^n\right)=\sum_n c(pn)q^n,
 \qquad
 \mathsf V_p\!\left(\sum_n c(n)q^n\right)=\sum_n c(n)q^{pn},
\]
and define
\begin{equation}
 B_U=2^{-\lceil\delta+w(U)\rceil}
             \prod_{p\in U}(\mathsf U_p-\eta(p)\mathsf V_p)g,
 \qquad
 e_U=\lceil\delta+w(U)\rceil-\delta-w(U).
 \label{tc:normalized-operator}
\end{equation}
The factors commute.  We say that $U$ has matching parity if
$\delta+w(U)$ is integral, and opposite parity otherwise.  Equivalently,
$w(U)$ has, or does not have, the fractional part prescribed by the filter.
Consequently $e_U=0$ or $1/2$, respectively.

These operations increase the level at a newly introduced $p$ only to
exponent one.  In fact, $\mathsf V_p$ changes the character by
$(p/\cdot)$ and raises the level by $p$, while on that raised level
$\mathsf U_p=\mathsf U_{p^2}\mathsf V_p$.  We denote the resulting
character of $B_U$ by $\psi_U$.  At a good prime $q\notin S\cup U$ the
coefficient recurrence remains
\begin{equation}
 B_U[q^2n]
 =\bigl(\eta(q)a_q(E)-\psi_U(q)(-n/q)\bigr)B_U[n]
       -q\eta(q)^2B_U[n/q^2].
 \label{tc:transformed-recurrence}
\end{equation}
For the odd source, Equation~\eqref{tc:transformed-recurrence} is read
modulo $2^{M-\lceil\delta+w(U)\rceil}$ when $M$ is the working
recurrence precision before normalization. As usual, a coefficient with
nonintegral index is zero. The equality of
the good eigenvalue before and after $\mathsf V_p$ can also be seen
directly: its new middle character contains $(p/q)$, which cancels the
factor $(p/q)$ arising from the change of coefficient index.

\begin{lemma}[Coefficient bound and first shell]
\label{tc:shell-bound}
Write an index prime to $S$ uniquely as
\[
             n=Jr m^2,
 \qquad J\subseteq U,
 \qquad r\text{ squarefree},\quad (r,U)=1,
\]
and put $D=(U/J)r$.

For the even source, $B_U$ has integral coefficients and satisfies the
following assertions at every index. For the odd source, fix bounds on
$w(U)$ and on the squarefree weights to be tested, and fix an integer
$K\ge0$. There is a larger cutoff $W$ and an ultrafilter-large set of
arbitrarily high stages at which $B_U$ is integral at every index, the
assertions below hold whenever $w(D)\le W$, and
\[
                  v(B_U[n])\ge K\qquad\text{if }w(D)>W.
\]
These choices are uniform over prime realizations with the prescribed
weight bounds.

The coefficient $B_U[n]$ is zero unless $D$ belongs to the filter.
Within the range just specified, it satisfies
\begin{equation}
 v(B_U[n])\geq
       w(r)+\frac{\#\{p\in J:p\text{ simple}\}}2-e_U.
 \label{tc:coefficient-bound}
\end{equation}
At a split prime $p\in J$, the factor contributed by the difference
operator, after removing its character phase and one factor of $2$, is
\begin{equation}
 c_p(D)=
 \frac{a_p(E)-(\psi/\eta)(p)(-D/p)-1}{2}\pmod 2.
 \label{tc:shell-factor}
\end{equation}
This is independent of the exponent of $p$ in $m$. At a simple prime
$p\in J$, Equation~\eqref{tc:shell-factor} holds for the first shell,
that is, when $p\nmid m$.
\end{lemma}

\begin{proof}
The good-prime recurrence of Section~\ref{cs:coefficient-systems} expresses
every relevant coefficient in terms of the squarefree coefficient $a(D)$.
After removing the factor $\eta(p)^s$, the multipliers for powers $p^{2s}$
are determined by
\[
 b_0=1,\qquad
 b_1=a_p(E)-(\psi/\eta)(p)(-D/p),\qquad
 b_{s+1}=a_p(E)b_s-pb_{s-1}.
\]
If $p\mid D$, the Legendre symbol in $b_1$ is zero.  The recurrences at
different primes commute and their multipliers multiply.

For $p\in U/J$ the difference operator contributes
$b_s-b_{s-1}$, with $b_{-1}=0$.  Here $p\mid D$, so $b_1=a_p(E)$ is
even.  Modulo $2$ the sequence alternates between $1$ and $0$, and every
such difference is odd.  For $p\in J$ we have $p\nmid D$; now $b_0$
and $b_1$ are odd and every $b_s$ is odd.  The difference
$b_{s+1}-b_s$ is therefore even.

For the more precise assertion at a split prime put
$a=a_p(E)$ and $\kappa=(\psi/\eta)(p)(-D/p)\in\{1,-1\}$.
Since $E[2]$ is rational over $\mathbf F_p$,
$a\equiv p+1\pmod4$.  If $\Delta_s=b_{s+1}-b_s$, the recurrence gives
\[
 \Delta_{s+1}=a\Delta_s-p\Delta_{s-1}\quad(s\geq1).
\]
Every $\Delta_s$ is even, whence
$\Delta_{s+1}\equiv\Delta_{s-1}\pmod4$.  Directly,
\[
 \Delta_1-\Delta_0
 =a^2-a\kappa-2a+2\kappa+1-p\equiv0\pmod4.
\]
Thus $\Delta_s/2\equiv(a-\kappa-1)/2\pmod2$ for every $s$.
For a simple prime this last formula is still the defining calculation
at $s=0$, which is all that will be required.

Put $d_U=\lceil\delta+w(U)\rceil$. Whenever the sharp squarefree
bound $v(a(D))\ge\delta+w(D)$ is available, the factors just computed
give at least one further factor of $2$ for each $p\in J$. Hence
\[
 v(B_U[n])\ge \delta+w(D)+\#J-d_U
 =w(r)+\frac{\#\{p\in J:p\text{ simple}\}}2-e_U.
\]
For an allowed squarefree class this expression is integral, since
$\delta+w(D)$ is integral. Its only possible negative value before
this parity condition is imposed is $-1/2$, so it is nonnegative.
For the even source the sharp bound holds at every squarefree index,
which proves all-index integrality and the displayed estimate.

For the odd source, let $u_0$ bound $w(U)$. Every prime in $U$ is a
nonroot, so $\#U\le2u_0$ and the normalizing exponents $d_U$ are
bounded. Proposition~\ref{prop:odd-coefficients} gives a constant
$C$ such that at every admissible squarefree index
\[
                   v(a_M(D))\ge\min\{M,w(D)-C\}.
\]
This includes the case of zero product derivative, when the coefficient
vanishes modulo $2^M$. Choose $W$ above the prescribed tested weights
and so large that $W-C-d_U>K$ for all permitted $U$. By
Lemma~\ref{cs:fixed-weight-uniformity}, on an ultrafilter-large set
of stages the sharp squarefree bound holds simultaneously for every
$D$ with $w(D)\le W$. Take an arbitrarily high stage in this set,
with precision exceeding all the fixed normalizations, the first-shell
tests, and $\delta+W+2u_0$, and such that $M-d_U\ge K$
for every permitted $U$.

The recurrences and operator identities for this source hold modulo
the working precision. Their multipliers are integral, so iteration
through any square exponent leaves errors in the same precision ideal.
After dividing by $2^{d_U}$, the possible error has valuation at least
$M-d_U$. Thus the sharp calculation above is valid throughout the
bounded range. Outside that range the same calculation with the coarse
squarefree bound gives
\[
 v(B_U[n])\ge
 \min\{M-d_U,\,w(D)-C+\#J-d_U\}\ge K.
\]
This proves integrality at every index of the actual form, while
asserting the sharp estimate only on the specified bounded range.
There is no restriction on the number of weight-zero root factors.
Coefficients excluded by the filter are zero. Every removed phase
is a value of $\eta$, hence a unit reducing to $1$ at the chosen
prime above $2$.
\end{proof}

At equality, the normalized leading coefficient is $A(D)$ times the
displayed shell factors. In particular it is zero when $A(D)=0$.
For the odd source this assertion is used for the finite stabilized
symbol evaluations realized at the chosen stage. All omitted-prime
comparisons use the same source $g$ and minimum $\delta$.

In later coefficientwise divisions, choose $K$ to cover the largest
additional division and the desired final divisibility. If a least-weight
argument also requires all lower-weight symbols to vanish, that condition
holds simultaneously on an ultrafilter-large set: a failure would select
a bounded-weight sequence contradicting the stabilized minimum. Intersect
that set with the one in the lemma. The bounded range then supplies the
sharp calculation and any required zero symbols; the coefficients beyond
the cutoff already have all the necessary factors of $2$. Consequently
all coefficientwise divisions and vanishings used below hold on actual
integral forms at one chosen stage for each finite collection of tests.
No stage is required to work for every unbounded configuration.

\subsection{Removing unary coefficients}

The coefficient bound identifies the terms that can survive at the
lowest valuation. For matching $U$ in the rational-torsion case, these
lie on finitely many square classes. Subtracting the corresponding unary
theta series will permit one further integral division by $2$.

Suppose first that $E(\Q)[2]\ne0$.  There are then no root primes.
For matching $U$, define
\begin{equation}
 \begin{split}
 \Theta_J&=
    \sum_{\substack{m>0\\(m,\prod_{s\in S}s)=1}}
       \eta(m)\chi_{-4}(m)^i m\,q^{Jm^2},\\
 H_U&=\frac12\left(B_U-
       \sum_{\substack{J\subseteq U\\J\text{ split only}}}
                    B_U[J]\Theta_J\right),
 \end{split}
 \label{tc:unary-subtraction}
\end{equation}
where $i\in\{0,1\}$ is chosen so that $\eta\chi_{-4}^i$ is odd.
For opposite $U$, set $H_U=B_U$.  Except where a further division is
explicitly made in Lemma~\ref{tc:minimum-transfer}, in the irreducible
case we also set $H_U=B_U$.

\begin{lemma}
\label{tc:unary-integrality}
The form in Equation~\eqref{tc:unary-subtraction} is integral.  Each
$\Theta_J$ is a weight-$3/2$ unary theta form with character
\[
                \eta\chi_{-4}^i(-J/\cdot),
\]
whose new level primes are the primes of $J$, each to exponent one.
Its character differs from $\psi_U$ by a quadratic character.
\end{lemma}

\begin{proof}
In matching parity the right side of
Equation~\eqref{tc:coefficient-bound} can be zero only if $r=1$ and
$J$ consists of split primes.  At these indices
Equation~\eqref{tc:shell-factor}, together with the odd squared-index
multipliers at primes outside $U$, gives
\[
                  B_U[Jm^2]\equiv B_U[J]\eta(m)\pmod2.
\]
The coefficient of $\Theta_J$ at $Jm^2$ has precisely this residue:
the common phase $\eta(m)$ is retained on both sides, while
$\chi_{-4}(m)^i m$ is an odd integer.  Thus this is divisibility by
$2$, not merely by $\mathfrak m$ in a ramified coefficient ring.
Different squarefree $J$ have disjoint such supports.  The numerator in
Equation~\eqref{tc:unary-subtraction} is consequently divisible by $2$
coefficientwise.

Put $\beta=\eta\chi_{-4}^i$, and let $\beta_0$ be its primitive
inducing character, of conductor $r$ supported on $S$.  The unary
theta transformation \citep[Section~2, p.~2]{Purkait2012} gives
\[
 \Theta_J^0=\sum_{m>0}\beta_0(m)m q^{Jm^2}
 \in S_{3/2}\bigl(4r^2J,\beta_0(-4J/\cdot)\bigr).
\]
Removing the terms with $s\mid m$ applies
$1-s\beta_0(s)\mathsf V_{s^2}$, where
$\mathsf V_{s^2}=\mathsf V_s^2$.  Thus
\[
 \Theta_J=\prod_{s\in S}
       \bigl(1-s\beta_0(s)\mathsf V_{s^2}\bigr)\Theta_J^0.
\]
Writing $P_S=\prod_{s\in S}s$, every term in this product has theta
parameter $Jd^2$ for some $d\mid P_S$.  The same transformation formula
places all terms at level $4r^2JP_S^2$, with character
$\beta_0(-4J/\cdot)$ on the units of that level.  This is the displayed
character $\eta\chi_{-4}^i(-J/\cdot)$.  Since $J$ is prime to $S$,
each fresh prime in $J$ still occurs in the level to exponent one.
Finally $\psi_U/\eta$ is quadratic, as is
$\chi_{-4}^i(-J/\cdot)$, proving the last assertion.
\end{proof}

\subsection{Weight-two traces}

We next express coefficient tests as traces of Galois representations.
The trace identities will give relations among the tests; the
level-prime identity will compare a test with the one obtained by
omitting a prime.

Multiply the integral half-integral-weight form under consideration by
\[
                  \theta=1+2\sum_{m\geq1}q^{m^2}.
\]
This produces a weight-two modular form with the same reduced expansion.
Choose a finite extension of $\Qtwo$ containing all coefficient fields
and character values needed at the fixed level, and denote its ring of
integers and maximal ideal by $\mathcal O$ and $\mathfrak m$.

The characteristic-zero decomposition by the Hecke operators away from
the level supplies commuting operators
\[
               t(u),\ d(u),\qquad u\in G_{\Q},
\]
by taking the traces and determinants of the corresponding
two-dimensional Galois representations.  For cuspidal systems these
are the weight-two representations of Deligne; for Eisenstein systems
they are the sums of the two defining characters.  At a good arithmetic
Frobenius, $t(\Fr_q)$ is the ordinary weight-two $T_q$, and the
determinant is the cyclotomic character times the diamond character.
We use the existence and good-prime compatibility of these representations
in their characteristic-zero form
\citep[\S3, (3.9)--(3.12), and Theorem~4.9]{Deligne1969}, and
local-global compatibility
\citep[Theorem~(A), \S0.7]{Carayol1986}.
We take the contragredients of the cohomological representations in these
references, so arithmetic Frobenius has trace $a_q$ and determinant $q$
times the diamond eigenvalue.

Here is a precise integral framework for the reductions that follow.
Take the finite-dimensional sum of these characteristic-zero systems,
retaining their multiplicity spaces, that contains the vector under
consideration and its diamond transforms.  Its integral Fourier lattice
is finite over $\mathcal O$: finitely many coefficients already give an
injective map of this space, and the lattice embeds in the corresponding
finite free $\mathcal O$-module.  Good Hecke operators preserve integral
expansions.  Diamonds do so as well.  For the subtracted vector this last
claim follows from the formula
\[
 \langle u\rangle H_U
 =\psi_U(u)H_U+
 \frac12\sum_J B_U[J]
       \bigl(\psi_U(u)-\eta(u)\chi_{-4}(u)^i(-J/u)\bigr)\Theta_J:
\]
each parenthesis is zero or twice a unit.  The same observation applies
after multiplying by $\theta$.

Traces at arbitrary Galois elements are limits of good Frobenius traces
on this finite lattice.  Indeed, continuity gives a common finite
quotient at any prescribed accuracy, and Chebotarev gives a good prime
in each required conjugacy class.  Its size is unrestricted, and it may
avoid any specified finite list of indices and primes.  This proves
integrality of all the traces and determinants used below, not only of
the good Hecke operators with which the construction began.

Let $\mathcal T$ be the resulting commutative operator algebra on the
reduced orbit of the vector. Equality means equality of every Fourier
coefficient after any further indicated operator products. We retain
this algebra, including its possible nilpotents.

\begin{lemma}[Trace identities and fresh-prime tests]
\label{tc:trace-algebra}\label{co:traces}
In this reduced operator algebra,
\begin{equation}
 d(u)t(v)=t(v),\qquad
 t(u)t(v)=t(uv)+t(uv^{-1}),\qquad t(1)=0.
 \label{tc:trace-identities}
\end{equation}
A coefficient of a product of traces can be tested by replacing its
arguments with Frobenius elements at distinct fresh primes avoiding
the coefficient index.  If the index is $n$ and those primes are
$q_1,\ldots,q_a$, the test is the reduced coefficient at
$nq_1\cdots q_a$.  Determinants may be applied first.
\end{lemma}

\begin{proof}
For a weight-two form $f$, the good-prime formula is
\[
                  (T_qf)[n]=f[qn]+q(\langle q\rangle f)[n/q].
\]
When $q\nmid n$ the second term is absent.  Choosing the successive
test primes distinct and avoiding $n$ proves the coefficient assertion.
Chebotarev and continuity, on the finite lattice just described, allow
these choices for arbitrary trace arguments at any required accuracy.

All the diamond characters are of $2$-power order.  On a form with one
character its determinant therefore reduces to the identity.  For a
subtracted vector the displayed diamond formula leaves, in addition,
only unary contributions $\Theta_J\theta$.  A fresh trace kills each
such contribution modulo $\mathfrak m$: its coefficient test has
index $nq$ with $q\nmid nJ$, and cannot have the squarefree part $J$.
The same argument applies after any further tests.  This proves
$d(u)t(v)=t(v)$.

For two invertible $2$-by-$2$ matrices the characteristic-zero identity
is
\[
 \tr(A)\tr(B)=\tr(AB)+\det(B)\tr(AB^{-1}).
\]
Apply it to every characteristic-zero system and use the determinant
identity just proved.  Finally, $t(1)=2$ before reduction.
\end{proof}

The bounds of Lemma~\ref{tc:shell-bound} are particularly well suited to
these tests: an added prime contributes its nonnegative weight, and
additional root primes contribute no weight.  A zero established by
one of these bounds consequently remains zero after all the further
tests used here.  In the moving-precision case all the preceding
constructions are made at the actual fixed level of the chosen stage.
The required divisibility is uniform at the specified weight cutoff;
continuity is used only in a finite-dimensional space at that stage.

The only level-prime identity we need is the following.  We distinguish
the weight-two operator $\mathsf U_p^{(2)}$ from the earlier
half-integral-weight notation, although both extract $p$-multiples in
their respective Fourier expansions.

\begin{lemma}[Inertia and a level operator]
\label{tc:inertia-identity}\label{co:inertia}
For $p\in U$, choose a place over $p$, an inertia element $\tau_p$
on which the ramified quadratic character is $-1$, and a corresponding
arithmetic Frobenius lift $\phi_p$.  On the characteristic-zero space
under consideration, before reduction, one has
\begin{equation}
 \mathsf U_p^{(2)}
       \bigl(t(g\tau_p\phi_p)-t(g\phi_p)\bigr)
 =d(\phi_p)\bigl(t(g\tau_p)-t(g)\bigr).
 \label{tc:exact-level-identity}
\end{equation}
The identity commutes with further good traces.  After reduction it
therefore evaluates a difference at coefficient $1$ by the corresponding
difference at coefficient $p$.
\end{lemma}

\begin{proof}
At $p$ the summands in question are either old from a level omitting
$p$, or have conductor exponent one and primitive ramified quadratic
diamond character.  For the former the representation is unramified,
so both trace differences are zero, irrespective of the action of
$\mathsf U_p^{(2)}$ on the old multiplicity space.

For the latter, local-global compatibility identifies the local
representation with a principal series with one unramified and one
ramified character; its monodromy is zero.  In a basis of these two
lines write
\[
 \rho(\tau_p)=\begin{pmatrix}1&0\\0&-1\end{pmatrix},
 \qquad
 \rho(\phi_p)=\begin{pmatrix}\lambda&0\\0&\mu\end{pmatrix}.
\]
The level eigenvalue is $\lambda$, the Frobenius eigenvalue on the
unramified line, and $d(\phi_p)=\lambda\mu$.  If the lower-right entry
of $\rho(g)$ is $z$, the two trace differences are $-2z\mu$ and
$-2z$, respectively.  Both sides of
Equation~\eqref{tc:exact-level-identity} are $-2z\lambda\mu$.
The Eisenstein calculation is the identical calculation on its pair
of characters.  Copies old at other primes do not change it.

These are all the possibilities: the new level has exponent one at
$p$, with ramified quadratic determinant, so the Steinberg case with
unramified determinant does not enter.  A unary summand with $p\notin J$
is old from a level omitting $p$; with $p\in J$ it has the same ramified
quadratic local character.  Thus the preceding classification also
covers Equation~\eqref{tc:unary-subtraction}.

Finally apply Equation~\eqref{tc:trace-identities} to remove the
determinant from the right side after reduction.  Taking coefficient
$1$ on the left is taking coefficient $p$ before $\mathsf U_p^{(2)}$.
\end{proof}

\subsection{Rules when there is rational two-torsion}

We now apply the common trace identities to the rational-torsion case.
The coefficient bounds first restrict the Galois data on which a symbol
can depend. The level-prime identity then determines how the symbol
changes when one of those data is varied.

An elementary extension of a number field means a compositum of
quadratic extensions.  Whenever a ramification support is specified,
only extensions unramified outside that support are included.  For
finite support the resulting elementary extension is finite.  Labels
below are Frobenius specifications in these extensions, with a chosen
place when an orientation over $F$ is needed.  They always include the
fixed rational local conditions that specify the filter.

\begin{proposition}[Rational trace and contraction rules]
\label{tc:rational-rules}
Suppose $E(\Q)[2]\ne0$.
\begin{enumerate}[label=\textup{(\roman*)}]
\item For matching $U$, on the subtracted vector $H_U$,
\[
                         t(g)t(u)=0\quad(u\text{ split}).
\]
The trace on split elements factors through the elementary extension
of $F$ supported on $S\cup U$.  For opposite $U$, on $B_U$, every
product of two traces is zero, and traces factor through the rational
elementary extension supported there.
\item If $U\cup\{q\}$ has the parity of the filter, its symbol is
read at coefficient $1$ by $t(\Fr_q)$ on $H_U$ for split $q$, and on
$B_U$ for simple $q$.
\item Fix $p\in U$ and put $g=\Fr_q$.  Replacing $g$ by $g\tau_p$
changes the symbol, additively in $\mathbf F_2$, by
\begin{equation}
                  A\bigl((U/p)\cup\{q'\}\bigr),
       \qquad \Fr_{q'}\sim g\phi_p,
 \label{tc:rational-contraction}
\end{equation}
if $p$ and $q$ are not both simple.  The label of $q'$ is required
only away from the omitted primes and carries their product labels.
The replacement uses the original source with level support
$S\cup(U/p)$: the \emph{whole rational prime} $p$ is omitted,
including both places over it if $p$ splits in $F$. If both endpoints
are simple, the change is zero.
\item A split vertex with trivial label relative to the level with
that vertex omitted has symbol zero.
\end{enumerate}
\end{proposition}

\begin{proof}
For matching $U$, two fresh tests, one of them split, add weight at
least $3/2$ to Equation~\eqref{tc:coefficient-bound}.  The allowed
valuation is integral, so the numerator has valuation at least $2$.
After the division in Equation~\eqref{tc:unary-subtraction} it still
vanishes modulo $\mathfrak m$; fresh tests kill its unary terms.
For opposite $U$, two nonroot tests add weight at least $1$, leaving
at least $1/2$ after subtracting $e_U$, and hence at least $1$ at an
allowed coefficient.  This proves the two product vanishings on the
entire reduced orbit.

The trace identities show why these vanishings give the stated
elementary quotients.  If $u$ is split, apply the identity to
$t(gu)t(u)$ to obtain $t(gu^2)=t(g)$.  The subgroup generated by
squares in $G_F$ has elementary quotient, since a group of exponent
two is abelian.  In the opposite case the same calculation applies
to every $u\in G_\Q$.  Ramification outside the level is absent from
the characteristic-zero representations and remains absent in the
quotient.

For a split $q$, its added weight is $1$: at coefficient $q$ the
division by $2$ in $H_U$ places the coefficient exactly at the
normalized detection depth.  There is no unary term at that index.
For simple $q$, $e_U=1/2$ already gives that depth on $B_U$.
The character phases reduce to $1$, proving (ii).

For (iii), apply Lemma~\ref{tc:inertia-identity} at coefficient $1$,
so that its left side is a coefficient-$p$ test. Replace its trace
arguments by fresh primes with Frobenius $g\tau_p\phi_p$ and
$g\phi_p$. Put $V=U/p$. For $X=U,V$, set
$n_X=\lceil\delta+w(X)\rceil$ and let $m_X$ be one in matching
parity and zero otherwise. After a fresh
trace kills the unary terms, the shell multiplier relative to the
omitted-source test is
\[
 2^{1+n_V-n_U-m_U+m_V}c_p(D'),
 \qquad D'=Vq',\quad \Fr_{q'}\sim g\phi_p.
\]
The exponent is zero for split $p$, and also for simple $p$ with
matching $U$. It is one for simple $p$ with opposite $U$, precisely
the simple--simple case. This checks all integer divisions in the
contraction rule. For simple $p$ only the first shell is used.
The change $\psi_U/\psi_V=(p/\cdot)$ is already included in that
shell calculation; all remaining phases reduce to one.

In the first three cases the shell is therefore $c_p(D')A(D')$.
Both fresh Frobenius elements have the same image in the omitted
source: that source is built from the original modular form, the same
$\eta,\psi,\delta$, and the operators at $V$, so it is unramified at the entire rational prime
$p$. Inertia at either place over $p$ acts trivially on its symbol.
The ramified quadratic coordinate of $\tau_p$ reverses the Legendre
sign in $c_p$, whose additive change is one. The difference of the
two shell values is consequently $A(D')$, proving
Equation~\eqref{tc:rational-contraction}.

If both primes are simple, their residual images are the same
transposition, so $g\phi_p$ is split.  The shell now has positive
valuation after normalization and contributes zero.  Finally (iv)
is (ii) with $t(\Fr_q)=t(1)=0$ in the relevant elementary quotient.
\end{proof}

The contraction rules control changes in existing vertices. The next
calculation gives a way to add two simple vertices while preserving the
symbol.

\begin{proposition}[Insertion of a pair of simple primes]
\label{tc:simple-twins}
Suppose $E(\Q)[2]\ne0$ and $U$ has matching parity.  Let a rational
elementary label, relative to $S$ and all existing split vertices,
be realized by a simple prime $q$ with
\[
                         \#E(\mathbf F_q)\equiv2\pmod4.
\]
Inserting two fresh simple primes with this same label leaves $A(U)$
unchanged.  Their mutual simple-edge label, and their labels against
other simple vertices, do not affect this assertion.
\end{proposition}

\begin{proof}
First test $t(\Fr_q)^2$ on matching $H_U$.  The weight-two Hecke
formula at coefficient $1$ gives, after reduction,
\[
                    H_U[q^2]+q(\langle q\rangle H_U)[1].
\]
Only the $J=1$ unary term can contribute at $1$ or $q^2$.
Write $b=B_U[1]$, $\epsilon=\chi_{-4}(q)$, and
$\kappa=(\psi_U/\eta)(q)$.  Substituting
Equation~\eqref{tc:transformed-recurrence} and the two characters in
Equation~\eqref{tc:unary-subtraction}, the preceding expression is
the reduction of
\begin{equation}
 b\,\frac{\eta(q)}2
 \left(a_q(E)+\kappa(q-\epsilon)
                 -q\epsilon^i(1+\epsilon)\right).
 \label{tc:twin-calculation}
\end{equation}
For $q\equiv3\pmod4$ the parenthesis divided by $2$ is
$a_q(E)/2$ modulo $2$.  For $q\equiv1\pmod4$ it is
$a_q(E)/2+1$.  In either case it equals
$(q+1-a_q(E))/2=\#E(\mathbf F_q)/2$ modulo $2$.
Since $b$ reduces to $A(U)$, this proves that the squared trace
has coefficient $(\#E(\mathbf F_q)/2)A(U)$.

By Lemma~\ref{tc:trace-algebra}, replace the two copies of this
trace by distinct fresh primes realizing the same operator tests.
Their product coefficient reads the symbol with the two primes
inserted.  To pass to any other pair with the labels in the
statement, vary one prime at a time.  After omitting that simple
prime the opposite-parity test of
Proposition~\ref{tc:rational-rules} depends only on rational
elementary labels.  Its changes against other simple primes are
zero by the simple--simple case of the same proposition.  Thus
only the fixed support and existing split labels need be matched.
Fresh intermediate primes can match both old and new finite
specifications, so that no collision of vertices is required.
The assumed odd value of $\#E(\mathbf F_q)/2$ completes the proof.
\end{proof}

\subsection{An elementary character on root elements}

Assume now that $E[2]$ is irreducible.  Its image is either cyclic
of order three or $S_3$.  The reduced trace algebra need not be a
field: the order-two units below can have nontrivial nilpotent
parts.  We therefore retain the algebra, rather than replacing it
by its residue field.

\begin{lemma}[Root character]
\label{tc:root-character}
On matching $B_U$ one has $t(u)=0$ for every nonroot $u$.
On opposite $B_U$ one has $t(u)=0$ for split $u$ and
$t(g)t(u)=0$ for two simple elements.  In either parity,
\[
                         t(u)^2=1\quad(u\text{ a root}).
\]
The traces factor through an extension of $\Gal(F/\Q)$ by an
elementary kernel over $F$.  Moreover, there is a multiplicative
character $\alpha:G_\Q\longrightarrow\mathcal T^\times$ into the
reduced trace algebra, unramified outside $S\cup U$, such that
\[
                  \alpha(u)^2=1\quad\text{for all }u,
 \qquad       t(u)=\alpha(u)\quad(u\text{ a root}).
\]
In particular $\alpha$ factors through a rational elementary
extension with that ramification support.
\end{lemma}

\begin{proof}
The trace vanishings follow from
Equation~\eqref{tc:coefficient-bound}.  In matching parity, one
nonroot test adds at least $1/2$ and hence gives an integral
valuation at least $1$.  In opposite parity, either a split test
or two simple tests add at least $1$, again leaving positive
integral valuation after the loss $e_U=1/2$.

For a root test, use the same fresh prime $q$ twice, avoiding the
coefficient index, rather than two distinct primes.  The weight-two
Hecke formula and Equation~\eqref{tc:transformed-recurrence} give
the identity operator modulo $\mathfrak m$: $a_q(E)$ is odd, and
the middle quadratic-character term cancels the diamond term
modulo $2$.  This proves $t(u)^2=1$ by continuity for every root
$u$.  For completeness, it suffices to check indices prime to the
chosen $q$; each particular index permits such a choice of $q$
with the prescribed finite trace data.

Since split traces vanish in both parities,
$t(gu^2)=t(g)$ for $u\in G_F$.  Thus the traces factor through
the quotient by the squares of $G_F$.  At any fixed accuracy the
trace action is on a finite lattice modulo $\mathfrak m$, so we
may work in a finite quotient
\[
       1\longrightarrow H\longrightarrow\Gamma
          \longrightarrow\Gal(F/\Q)\longrightarrow1
\]
with $H$ elementary abelian.  Choose an order-three lift $R$ of a
rotation.  Such a lift exists because the inverse image of the
order-three subgroup is an extension by a finite $2$-group.
The trace identity with the identity element gives
$t(R)=t(R^{-1})=t(R^2)$.  Applying it to $R,R^2$ gives
$t(R)t(R^2)=t(1)+t(R^{-1})=t(R)$; the left side is $1$.
Hence
\[
                           t(R)=t(R^2)=1.
\]

For $a\in H$ set
\[
                       \alpha(a)=t(aR)=t(aR^2).
\]
The equality follows from $t(a)=0$ and the trace identity with
$R$.  Conjugacy invariance by an element $b\in H$ implies that
$\alpha$ is unchanged on multiplying its argument by
$bR(b)^{-1}$.  It therefore factors through the coinvariants
$H/[H,R]$.  Also it is invariant under conjugation by $R$.
The trace product now gives
\[
 \alpha(a)\alpha(b)
   =t(aR)t(bR)
   =t\bigl(aR(b)R^2\bigr)+t(ab^{-1})
   =\alpha\bigl(aR(b)\bigr)
   =\alpha(ab).
\]
The last equality uses the coinvariant invariance just proved;
the split term is zero.  Since $aR$ is a root,
$\alpha(a)^2=1$.

For cyclic residual image, this already extends to $\Gamma$ by
setting $\alpha(R)=1$.  For $S_3$, choose a transposition lift
$s$ normalizing $\langle R\rangle$.  It can be obtained from any
transposition lift by conjugating its order-three complement by
an element of $H$.  Thus $sRs^{-1}=R^{-1}$ and $s^2\in H$.
Conjugacy of traces shows that $\alpha$ on $H$ is invariant
under $s$.  In both parities we have $t(s)t(sR)=0$, either because
each factor is zero or by the two-simple product rule.  Its trace
identity reads
\[
                       0=\alpha(s^2)+t(R)=\alpha(s^2)+1.
\]
Therefore the extension defined by $\alpha(s)=1$ respects the
relation for $s^2$, as well as conjugation and the relations in
$\langle R\rangle$.  It is the required multiplicative
order-two character, and agrees with the trace on both root
cosets.  At the fixed stage, choose one common finite quotient
through which the traces and residual representation factor, and
pull back the character constructed there.  No canonical choice
across quotients is needed.  Unramifiedness outside the level is preserved.
\end{proof}

\subsection{Irreducible deletion rules and minimal labels}

The root character replaces products of root traces by a single
multiplicative value. We now combine that fact with the level-prime
identity to compute changes at nonroot vertices. The last lemma applies
these rules at the least weight where a symbol is nonzero.

The \emph{total root label} is the product of the rational elementary
labels of all root primes being used.  It includes the fixed-support
conditions and the rational quadratic labels against the active
nonroot vertices.  Its discriminant coordinate is trivial in the
quadratic quotient of $S_3$; for cyclic residual image there is no
such coordinate.

\begin{proposition}[Irreducible substitution rules]
\label{tc:irreducible-rules}
Suppose $E[2]$ is irreducible.
\begin{enumerate}[label=\textup{(\roman*)}]
\item On matching $B_U$, coefficient $1$ of a product of root
traces reads $A$ on the union of $U$ and those roots.  It depends
on the roots only through their total rational elementary label
relative to $S\cup U$.  Any finite collection of roots, including
the empty collection, may be replaced in this test by one virtual
root with that label.
\item Changing the Legendre bit of that root against a split
vertex $p$ has additive derivative equal to deletion of $p$,
with its rational labels away from $p$ transferred to the total
root label.  Changing the root bit against a simple vertex has
derivative zero.
\item To vary a simple vertex $q$, omit it and use opposite
$B_V$, retaining the other root probes.  At another simple
vertex $p$, choose Frobenius lifts $g$ and $\phi_p$ for $q$ and
$p$.  If their transposition images are different, the derivative
under $g\mapsto g\tau_p$ deletes both $p$ and $q$ and transfers
their rational labels to the root total.  If their transposition
images are equal, the derivative is zero.
\item In the test in (iii), variation at a split $p$ is given by
the first-shell factors $c_p$ times symbols of merged
configurations whose weight is one less.  It consequently
vanishes whenever all symbols of that lower weight vanish.
\end{enumerate}
All these assertions preserve every specified label at the other
vertices, and use the chosen place for inertia and Frobenius.
\end{proposition}

\begin{proof}
Root primes have weight zero.  Thus the fresh-prime tests at
coefficient $1$ of matching $B_U$ are precisely the normalized
coefficients defining $A$.  Lemma~\ref{tc:root-character} makes
the product of root traces equal to the value of $\alpha$ on
their total elementary label.  A rotation survives every
compatible elementary specification: an elementary quotient
cannot eliminate the order-three subgroup, and in the $S_3$
case the only compatibility condition is its trivial image in
the quadratic quotient.  Chebotarev therefore supplies the
claimed single virtual root, even for the empty total label.

For (ii), apply Lemma~\ref{tc:inertia-identity} with $g$ a root.
If $p$ is split then $g\phi_p$ is a root.  The coefficient with
$J=\{p\}$ has zero cost in
Equation~\eqref{tc:coefficient-bound}; its first shell is
$c_p$ times the symbol with $p$ deleted.  The latter is
unramified at $p$ and is the same in the two tests, while $c_p$
changes by $1$.  Multiplication of Frobenius labels by $\phi_p$
is exactly the transfer of the deleted vertex's rational labels
to the root.  If $p$ is simple, $g\phi_p$ is simple, and its
weight $1/2$ plus the $1/2$ contributed by $J$ puts the shell
above the testing depth.  The derivative is then zero.

For (iii), $e_V=1/2$.  When the two transpositions differ,
$g\phi_p$ is a rotation, and the $1/2$ contribution from the
simple prime in $J=\{p\}$ exactly cancels $e_V$.  The shell
again equals $c_p$ times the omitted-$p$ symbol, now with both
simple vertices removed and their product represented by a
root.  The Legendre flip proves the asserted derivative.
For equal transpositions their product is split; its additional
weight $1$ forces the shell to vanish.

Finally, if $p$ is split and $g$ simple, then $g\phi_p$ is
simple.  The $p$-shell is at the first-shell depth and is
expressed by Equation~\eqref{tc:shell-factor} with
$D'=(V/p)q'$ and the retained root probes.  Its nonroot weight
is
\[
                 w(V)-1+\frac12
                  =\left(w(V)+\frac12\right)-1.
\]
This is one less than the original weight.  In particular both
shell terms vanish when symbols of that weight vanish.  The
comparison omitting $p$ is unramified there.  Root probes can
be matched on the rational elementary support throughout,
which proves the final label assertion as well.
\end{proof}

At the least nonzero weight, every lower-weight term in the preceding
rules vanishes. This removes the dependence on internal edge labels;
the remaining assertion keeps both the fixed-support labels and the
total root label.

\begin{lemma}[Transfer at the least nonzero weight]
\label{tc:minimum-transfer}
Suppose $E[2]$ is irreducible, and let $m$ be the least weight
of an index on the filter with nonzero symbol.  At weight $m$,
the dependence on each nonroot vertex is only through its
oriented rational Frobenius lift in the Galois group over $\Q$
of the maximal elementary extension of $F$ unramified outside $S$, together with the
requirement on the total root label.  In particular, replacing
the nonroot vertices by fresh vertices with those same
fixed-support labels preserves the symbol, irrespective of
their internal edge labels, provided the root total is matched.
The assertion allows simultaneous changes using disjoint
intermediate configurations.
\end{lemma}

\begin{proof}
Every symbol of weight less than $m$ is zero.  For a simple
vertex, Proposition~\ref{tc:irreducible-rules} therefore makes
its inertia derivatives at every other active nonroot vertex
zero: different transpositions give a lower-weight deletion,
equal transpositions give zero directly, and a split vertex
gives a lower-weight merged configuration.  Its labels against
root vertices are immaterial by the simple case of the root
derivative rule.  These statements hold for conjugates of
inertia too, because the Frobenius and inertia calculation is
made at a chosen place and may be conjugated.  Starting with
the elementary extension over $F$ supported on $S\cup V$,
killing the inertia at the active vertices leaves exactly
fixed-support elementary labels.  This proves the assertion
for a simple vertex.

For a split vertex $q$, let $V$ be the other nonroot vertices,
so $w(V)=m-1$ and $V$ still has matching parity.  We first claim
that $B_V$ is divisible by $2$ integrally.  In
Equation~\eqref{tc:coefficient-bound}, an entry of cost zero
has $r$ consisting only of roots and $J$ consisting only of
split primes.  Its underlying squarefree class has weight
\[
                         w(V)-\#J<m.
\]
Its symbol is zero.  The shell formula at every split prime,
valid at all exponents, therefore adds one more factor of $2$
to the normalized coefficient.  All other allowed entries
already have positive integral cost.  This proves the claim,
including at arbitrarily large coefficient indices.

Use now $H_V=B_V/2$.  Two split traces give zero on $H_V$:
their added weight $2$ in
Equation~\eqref{tc:coefficient-bound} still leaves one factor
of $2$ after this division.  The trace restricted to split
elements consequently factors through the elementary kernel
over $F$, by the same square-invariance argument as before.
Coefficient $1$ with the new split trace and the retained
root probes reads the symbol at weight $m$.

Apply Lemma~\ref{tc:inertia-identity} at $p\in V$ to this
divided vector.  The element $g\phi_p$ has the type of $p$,
because $g$ is split.  Its $p$-shell is $c_p$ times a symbol
of weight $m-1$, whether $p$ is split or simple; for a simple
$p$ only the first shell is used.  For a simple $p$ the normalization can be checked directly:
$J=\{p\}$ and the new simple probe contribute $1/2+1/2=1$ to
Equation~\eqref{tc:coefficient-bound}. Writing
$n_V=\delta+w(V)$, the divided shell is
\[
 \frac{2c_p(D')a(D')}{2^{n_V+1}}
 =c_p(D')\frac{a(D')}{2^{n_V}},\qquad w(D')=w(V)=m-1.
\]
The omitted set $V/p$ has opposite parity, but its ceiling
normalization is still $n_V$. Thus this is exactly its unit-symbol
test, with no fractional loss. Both compared values are zero.
Every such inertia derivative therefore vanishes, leaving only the
fixed-support elementary label.

No character $\alpha$ on this further divided vector is
needed.  Use actual matched root probes in both tests.  At
weight $m$, their changes against a split vertex are already
zero by the split root-deletion rule, because its deletion
has weight $m-1$.  They can therefore be matched while the
nonroot vertex is moved.  Moving vertices successively,
using fresh intermediate copies when necessary, proves the
joint assertion without imposing old internal edge labels.

For a moving coefficient system, minimality is used on a fixed bounded
weight range. A failure of the required lower-weight vanishing on an
ultrafilter-large set of stages would give a bounded-weight sequence
contradicting the stabilized minimum. Choose the actual stage as in
Lemma~\ref{tc:shell-bound}, with the additional division by $2$ included
among its fixed division bounds. Its coarse estimate controls all larger
weights, and its integral square-index recurrences control every shell.
Thus the preceding coefficientwise arguments apply to the actual forms
at that stage. The argument does not require $m$ to be integral.
\end{proof}

\subsection{Summary and finite realization}

We collect the rules in the form used to construct finite prime arrays.
The analytic meaning of a nonzero symbol is unchanged; this section has
established which substitutions preserve it or compute its change.

\begin{theorem}[Symbol rules]
\label{thm:symbol-rules}\label{co:symbol-interface}
For either coefficient source of
Propositions~\ref{prop:even-coefficients} and
\ref{prop:odd-coefficients}, assume its stated coefficient
lower bound and use its normalized symbol and filter.
Then the following rules hold for finite configurations of
primes outside the level support.
\begin{enumerate}[label=\textup{(\roman*)}]
\item With rational two-torsion, the elementary-label
factorizations, contraction except in the simple--simple
case, and isolation by a trivial split label are those of
Proposition~\ref{tc:rational-rules}.  A pair of identical
simple labels realized with $\#E(\mathbf F_q)/2$ odd can be
inserted without changing the symbol, as in
Proposition~\ref{tc:simple-twins}.
\item With irreducible $E[2]$, all root probes combine into
one total rational elementary label.  A root-bit derivative
at a split vertex deletes that vertex, while at a simple
vertex it is zero.  The derivative between two simple
vertices deletes both when their transposition images are
different, and is zero when they are equal.  The remaining
simple--split shell has weight one less.  These are the
rules of Proposition~\ref{tc:irreducible-rules}.
\item At the least weight supporting a nonzero symbol in
the irreducible case, each nonroot may be moved with only
its fixed-support elementary Frobenius label prescribed;
internal edges impose no further condition on the value,
subject to matching the total root label.  This is
Lemma~\ref{tc:minimum-transfer}.
\end{enumerate}
The analytic and finite-Selmer consequences of a nonzero
symbol remain those of
Proposition~\ref{prop:unit-symbol-detection}.  In the odd
source the rules mean finite tests at arbitrarily high stages
chosen in the ultrafilter-large sets of the prescribed stabilization.
\end{theorem}

\begin{proof}
The preceding propositions prove the assertions. For one application,
put the finitely many Galois modules and operator tests in a common
finite quotient at the required precision. Chebotarev then gives
simultaneously fresh Frobenius representatives. The omitted source is
unramified at the whole omitted prime, whereas its ramified rational
quadratic coordinate is exactly the Legendre flip in the shell formula.
These are therefore substitutions with all other labels held fixed.

The arithmetic compatibility of an entire array is addressed in the
following sections. The present assertions concern each finite array;
in the odd case choose a sufficiently high actual stage in the
required ultrafilter-large intersection. Arbitrarily many weight-zero roots cause no additional depth
loss, because the coefficient bound is uniform at a fixed weight and
their traces combine through one total elementary label.

\end{proof}
